\documentclass[../../main.tex]{subfiles}
\begin{document}

\begin{flushright}
  \footnotesize\textit{【自動文字起こし・要確認】}
\end{flushright}

% Source: handwritten.pdf 第6問．
{\bf [解答]}
球および箱をすべて区別して考える．すべての場合の数は$(2n)^n$であり，どの箱にも球が2個以上入らない場合の数は
\begin{align}{}_{2n}P_n=(2n)(2n-1)\cdots(n+1).
\end{align}
したがって
\begin{align}
  p_n&=\prod_{k=n+1}^{2n}\dfrac{k}{2n},\\
  \log p_n&=\sum_{k=n+1}^{2n}\log\dfrac{k}{2n}.
\end{align}
区分求積法より
\begin{align}
  \lim_{n\to\infty}\dfrac{\log p_n}{n}
  &=2\int_{1/2}^1\log x\,dx\\
  &=2\left[x(\log x-1)\right]_{1/2}^1\\
  &=\log2-1.
\end{align}

\end{document}
